\documentclass[11pt]{article}

    \usepackage[breakable]{tcolorbox}
    \usepackage{parskip} % Stop auto-indenting (to mimic markdown behaviour)
    

    % Basic figure setup, for now with no caption control since it's done
    % automatically by Pandoc (which extracts ![](path) syntax from Markdown).
    \usepackage{graphicx}
    % Keep aspect ratio if custom image width or height is specified
    \setkeys{Gin}{keepaspectratio}
    % Maintain compatibility with old templates. Remove in nbconvert 6.0
    \let\Oldincludegraphics\includegraphics
    % Ensure that by default, figures have no caption (until we provide a
    % proper Figure object with a Caption API and a way to capture that
    % in the conversion process - todo).
    \usepackage{caption}
    \DeclareCaptionFormat{nocaption}{}
    \captionsetup{format=nocaption,aboveskip=0pt,belowskip=0pt}

    \usepackage{float}
    \floatplacement{figure}{H} % forces figures to be placed at the correct location
    \usepackage{xcolor} % Allow colors to be defined
    \usepackage{enumerate} % Needed for markdown enumerations to work
    \usepackage{geometry} % Used to adjust the document margins
    \usepackage{amsmath} % Equations
    \usepackage{amssymb} % Equations
    \usepackage{textcomp} % defines textquotesingle
    % Hack from http://tex.stackexchange.com/a/47451/13684:
    \AtBeginDocument{%
        \def\PYZsq{\textquotesingle}% Upright quotes in Pygmentized code
    }
    \usepackage{upquote} % Upright quotes for verbatim code
    \usepackage{eurosym} % defines \euro

    \usepackage{iftex}
    \ifPDFTeX
        \usepackage[T1]{fontenc}
        \IfFileExists{alphabeta.sty}{
              \usepackage{alphabeta}
          }{
              \usepackage[mathletters]{ucs}
              \usepackage[utf8x]{inputenc}
          }
    \else
        \usepackage{fontspec}
        \usepackage{unicode-math}
    \fi

    \usepackage{fancyvrb} % verbatim replacement that allows latex
    \usepackage{grffile} % extends the file name processing of package graphics
                         % to support a larger range
    \makeatletter % fix for old versions of grffile with XeLaTeX
    \@ifpackagelater{grffile}{2019/11/01}
    {
      % Do nothing on new versions
    }
    {
      \def\Gread@@xetex#1{%
        \IfFileExists{"\Gin@base".bb}%
        {\Gread@eps{\Gin@base.bb}}%
        {\Gread@@xetex@aux#1}%
      }
    }
    \makeatother
    \usepackage[Export]{adjustbox} % Used to constrain images to a maximum size
    \adjustboxset{max size={0.9\linewidth}{0.9\paperheight}}

    % The hyperref package gives us a pdf with properly built
    % internal navigation ('pdf bookmarks' for the table of contents,
    % internal cross-reference links, web links for URLs, etc.)
    \usepackage{hyperref}
    % The default LaTeX title has an obnoxious amount of whitespace. By default,
    % titling removes some of it. It also provides customization options.
    \usepackage{titling}
    \usepackage{longtable} % longtable support required by pandoc >1.10
    \usepackage{booktabs}  % table support for pandoc > 1.12.2
    \usepackage{array}     % table support for pandoc >= 2.11.3
    \usepackage{calc}      % table minipage width calculation for pandoc >= 2.11.1
    \usepackage[inline]{enumitem} % IRkernel/repr support (it uses the enumerate* environment)
    \usepackage[normalem]{ulem} % ulem is needed to support strikethroughs (\sout)
                                % normalem makes italics be italics, not underlines
    \usepackage{soul}      % strikethrough (\st) support for pandoc >= 3.0.0
    \usepackage{mathrsfs}
    

    
    % Colors for the hyperref package
    \definecolor{urlcolor}{rgb}{0,.145,.698}
    \definecolor{linkcolor}{rgb}{.71,0.21,0.01}
    \definecolor{citecolor}{rgb}{.12,.54,.11}

    % ANSI colors
    \definecolor{ansi-black}{HTML}{3E424D}
    \definecolor{ansi-black-intense}{HTML}{282C36}
    \definecolor{ansi-red}{HTML}{E75C58}
    \definecolor{ansi-red-intense}{HTML}{B22B31}
    \definecolor{ansi-green}{HTML}{00A250}
    \definecolor{ansi-green-intense}{HTML}{007427}
    \definecolor{ansi-yellow}{HTML}{DDB62B}
    \definecolor{ansi-yellow-intense}{HTML}{B27D12}
    \definecolor{ansi-blue}{HTML}{208FFB}
    \definecolor{ansi-blue-intense}{HTML}{0065CA}
    \definecolor{ansi-magenta}{HTML}{D160C4}
    \definecolor{ansi-magenta-intense}{HTML}{A03196}
    \definecolor{ansi-cyan}{HTML}{60C6C8}
    \definecolor{ansi-cyan-intense}{HTML}{258F8F}
    \definecolor{ansi-white}{HTML}{C5C1B4}
    \definecolor{ansi-white-intense}{HTML}{A1A6B2}
    \definecolor{ansi-default-inverse-fg}{HTML}{FFFFFF}
    \definecolor{ansi-default-inverse-bg}{HTML}{000000}

    % common color for the border for error outputs.
    \definecolor{outerrorbackground}{HTML}{FFDFDF}

    % commands and environments needed by pandoc snippets
    % extracted from the output of `pandoc -s`
    \providecommand{\tightlist}{%
      \setlength{\itemsep}{0pt}\setlength{\parskip}{0pt}}
    \DefineVerbatimEnvironment{Highlighting}{Verbatim}{commandchars=\\\{\}}
    % Add ',fontsize=\small' for more characters per line
    \newenvironment{Shaded}{}{}
    \newcommand{\KeywordTok}[1]{\textcolor[rgb]{0.00,0.44,0.13}{\textbf{{#1}}}}
    \newcommand{\DataTypeTok}[1]{\textcolor[rgb]{0.56,0.13,0.00}{{#1}}}
    \newcommand{\DecValTok}[1]{\textcolor[rgb]{0.25,0.63,0.44}{{#1}}}
    \newcommand{\BaseNTok}[1]{\textcolor[rgb]{0.25,0.63,0.44}{{#1}}}
    \newcommand{\FloatTok}[1]{\textcolor[rgb]{0.25,0.63,0.44}{{#1}}}
    \newcommand{\CharTok}[1]{\textcolor[rgb]{0.25,0.44,0.63}{{#1}}}
    \newcommand{\StringTok}[1]{\textcolor[rgb]{0.25,0.44,0.63}{{#1}}}
    \newcommand{\CommentTok}[1]{\textcolor[rgb]{0.38,0.63,0.69}{\textit{{#1}}}}
    \newcommand{\OtherTok}[1]{\textcolor[rgb]{0.00,0.44,0.13}{{#1}}}
    \newcommand{\AlertTok}[1]{\textcolor[rgb]{1.00,0.00,0.00}{\textbf{{#1}}}}
    \newcommand{\FunctionTok}[1]{\textcolor[rgb]{0.02,0.16,0.49}{{#1}}}
    \newcommand{\RegionMarkerTok}[1]{{#1}}
    \newcommand{\ErrorTok}[1]{\textcolor[rgb]{1.00,0.00,0.00}{\textbf{{#1}}}}
    \newcommand{\NormalTok}[1]{{#1}}

    % Additional commands for more recent versions of Pandoc
    \newcommand{\ConstantTok}[1]{\textcolor[rgb]{0.53,0.00,0.00}{{#1}}}
    \newcommand{\SpecialCharTok}[1]{\textcolor[rgb]{0.25,0.44,0.63}{{#1}}}
    \newcommand{\VerbatimStringTok}[1]{\textcolor[rgb]{0.25,0.44,0.63}{{#1}}}
    \newcommand{\SpecialStringTok}[1]{\textcolor[rgb]{0.73,0.40,0.53}{{#1}}}
    \newcommand{\ImportTok}[1]{{#1}}
    \newcommand{\DocumentationTok}[1]{\textcolor[rgb]{0.73,0.13,0.13}{\textit{{#1}}}}
    \newcommand{\AnnotationTok}[1]{\textcolor[rgb]{0.38,0.63,0.69}{\textbf{\textit{{#1}}}}}
    \newcommand{\CommentVarTok}[1]{\textcolor[rgb]{0.38,0.63,0.69}{\textbf{\textit{{#1}}}}}
    \newcommand{\VariableTok}[1]{\textcolor[rgb]{0.10,0.09,0.49}{{#1}}}
    \newcommand{\ControlFlowTok}[1]{\textcolor[rgb]{0.00,0.44,0.13}{\textbf{{#1}}}}
    \newcommand{\OperatorTok}[1]{\textcolor[rgb]{0.40,0.40,0.40}{{#1}}}
    \newcommand{\BuiltInTok}[1]{{#1}}
    \newcommand{\ExtensionTok}[1]{{#1}}
    \newcommand{\PreprocessorTok}[1]{\textcolor[rgb]{0.74,0.48,0.00}{{#1}}}
    \newcommand{\AttributeTok}[1]{\textcolor[rgb]{0.49,0.56,0.16}{{#1}}}
    \newcommand{\InformationTok}[1]{\textcolor[rgb]{0.38,0.63,0.69}{\textbf{\textit{{#1}}}}}
    \newcommand{\WarningTok}[1]{\textcolor[rgb]{0.38,0.63,0.69}{\textbf{\textit{{#1}}}}}
    \makeatletter
    \newsavebox\pandoc@box
    \newcommand*\pandocbounded[1]{%
      \sbox\pandoc@box{#1}%
      % scaling factors for width and height
      \Gscale@div\@tempa\textheight{\dimexpr\ht\pandoc@box+\dp\pandoc@box\relax}%
      \Gscale@div\@tempb\linewidth{\wd\pandoc@box}%
      % select the smaller of both
      \ifdim\@tempb\p@<\@tempa\p@
        \let\@tempa\@tempb
      \fi
      % scaling accordingly (\@tempa < 1)
      \ifdim\@tempa\p@<\p@
        \scalebox{\@tempa}{\usebox\pandoc@box}%
      % scaling not needed, use as it is
      \else
        \usebox{\pandoc@box}%
      \fi
    }
    \makeatother

    % Define a nice break command that doesn't care if a line doesn't already
    % exist.
    \def\br{\hspace*{\fill} \\* }
    % Math Jax compatibility definitions
    \let\Oldtex\TeX
    \let\Oldlatex\LaTeX
    \renewcommand{\TeX}{\textrm{\Oldtex}}
    \renewcommand{\LaTeX}{\textrm{\Oldlatex}}
    % Document parameters
    % Document title

    \input{../../macros.tex}
    \title{Homework\_13}
    
    
    
    
    
    
    
% Pygments definitions
\makeatletter
\def\PY@reset{\let\PY@it=\relax \let\PY@bf=\relax%
    \let\PY@ul=\relax \let\PY@tc=\relax%
    \let\PY@bc=\relax \let\PY@ff=\relax}
\def\PY@tok#1{\csname PY@tok@#1\endcsname}
\def\PY@toks#1+{\ifx\relax#1\empty\else%
    \PY@tok{#1}\expandafter\PY@toks\fi}
\def\PY@do#1{\PY@bc{\PY@tc{\PY@ul{%
    \PY@it{\PY@bf{\PY@ff{#1}}}}}}}
\def\PY#1#2{\PY@reset\PY@toks#1+\relax+\PY@do{#2}}

\@namedef{PY@tok@w}{\def\PY@tc##1{\textcolor[rgb]{0.73,0.73,0.73}{##1}}}
\@namedef{PY@tok@c}{\let\PY@it=\textit\def\PY@tc##1{\textcolor[rgb]{0.24,0.48,0.48}{##1}}}
\@namedef{PY@tok@cp}{\def\PY@tc##1{\textcolor[rgb]{0.61,0.40,0.00}{##1}}}
\@namedef{PY@tok@k}{\let\PY@bf=\textbf\def\PY@tc##1{\textcolor[rgb]{0.00,0.50,0.00}{##1}}}
\@namedef{PY@tok@kp}{\def\PY@tc##1{\textcolor[rgb]{0.00,0.50,0.00}{##1}}}
\@namedef{PY@tok@kt}{\def\PY@tc##1{\textcolor[rgb]{0.69,0.00,0.25}{##1}}}
\@namedef{PY@tok@o}{\def\PY@tc##1{\textcolor[rgb]{0.40,0.40,0.40}{##1}}}
\@namedef{PY@tok@ow}{\let\PY@bf=\textbf\def\PY@tc##1{\textcolor[rgb]{0.67,0.13,1.00}{##1}}}
\@namedef{PY@tok@nb}{\def\PY@tc##1{\textcolor[rgb]{0.00,0.50,0.00}{##1}}}
\@namedef{PY@tok@nf}{\def\PY@tc##1{\textcolor[rgb]{0.00,0.00,1.00}{##1}}}
\@namedef{PY@tok@nc}{\let\PY@bf=\textbf\def\PY@tc##1{\textcolor[rgb]{0.00,0.00,1.00}{##1}}}
\@namedef{PY@tok@nn}{\let\PY@bf=\textbf\def\PY@tc##1{\textcolor[rgb]{0.00,0.00,1.00}{##1}}}
\@namedef{PY@tok@ne}{\let\PY@bf=\textbf\def\PY@tc##1{\textcolor[rgb]{0.80,0.25,0.22}{##1}}}
\@namedef{PY@tok@nv}{\def\PY@tc##1{\textcolor[rgb]{0.10,0.09,0.49}{##1}}}
\@namedef{PY@tok@no}{\def\PY@tc##1{\textcolor[rgb]{0.53,0.00,0.00}{##1}}}
\@namedef{PY@tok@nl}{\def\PY@tc##1{\textcolor[rgb]{0.46,0.46,0.00}{##1}}}
\@namedef{PY@tok@ni}{\let\PY@bf=\textbf\def\PY@tc##1{\textcolor[rgb]{0.44,0.44,0.44}{##1}}}
\@namedef{PY@tok@na}{\def\PY@tc##1{\textcolor[rgb]{0.41,0.47,0.13}{##1}}}
\@namedef{PY@tok@nt}{\let\PY@bf=\textbf\def\PY@tc##1{\textcolor[rgb]{0.00,0.50,0.00}{##1}}}
\@namedef{PY@tok@nd}{\def\PY@tc##1{\textcolor[rgb]{0.67,0.13,1.00}{##1}}}
\@namedef{PY@tok@s}{\def\PY@tc##1{\textcolor[rgb]{0.73,0.13,0.13}{##1}}}
\@namedef{PY@tok@sd}{\let\PY@it=\textit\def\PY@tc##1{\textcolor[rgb]{0.73,0.13,0.13}{##1}}}
\@namedef{PY@tok@si}{\let\PY@bf=\textbf\def\PY@tc##1{\textcolor[rgb]{0.64,0.35,0.47}{##1}}}
\@namedef{PY@tok@se}{\let\PY@bf=\textbf\def\PY@tc##1{\textcolor[rgb]{0.67,0.36,0.12}{##1}}}
\@namedef{PY@tok@sr}{\def\PY@tc##1{\textcolor[rgb]{0.64,0.35,0.47}{##1}}}
\@namedef{PY@tok@ss}{\def\PY@tc##1{\textcolor[rgb]{0.10,0.09,0.49}{##1}}}
\@namedef{PY@tok@sx}{\def\PY@tc##1{\textcolor[rgb]{0.00,0.50,0.00}{##1}}}
\@namedef{PY@tok@m}{\def\PY@tc##1{\textcolor[rgb]{0.40,0.40,0.40}{##1}}}
\@namedef{PY@tok@gh}{\let\PY@bf=\textbf\def\PY@tc##1{\textcolor[rgb]{0.00,0.00,0.50}{##1}}}
\@namedef{PY@tok@gu}{\let\PY@bf=\textbf\def\PY@tc##1{\textcolor[rgb]{0.50,0.00,0.50}{##1}}}
\@namedef{PY@tok@gd}{\def\PY@tc##1{\textcolor[rgb]{0.63,0.00,0.00}{##1}}}
\@namedef{PY@tok@gi}{\def\PY@tc##1{\textcolor[rgb]{0.00,0.52,0.00}{##1}}}
\@namedef{PY@tok@gr}{\def\PY@tc##1{\textcolor[rgb]{0.89,0.00,0.00}{##1}}}
\@namedef{PY@tok@ge}{\let\PY@it=\textit}
\@namedef{PY@tok@gs}{\let\PY@bf=\textbf}
\@namedef{PY@tok@ges}{\let\PY@bf=\textbf\let\PY@it=\textit}
\@namedef{PY@tok@gp}{\let\PY@bf=\textbf\def\PY@tc##1{\textcolor[rgb]{0.00,0.00,0.50}{##1}}}
\@namedef{PY@tok@go}{\def\PY@tc##1{\textcolor[rgb]{0.44,0.44,0.44}{##1}}}
\@namedef{PY@tok@gt}{\def\PY@tc##1{\textcolor[rgb]{0.00,0.27,0.87}{##1}}}
\@namedef{PY@tok@err}{\def\PY@bc##1{{\setlength{\fboxsep}{\string -\fboxrule}\fcolorbox[rgb]{1.00,0.00,0.00}{1,1,1}{\strut ##1}}}}
\@namedef{PY@tok@kc}{\let\PY@bf=\textbf\def\PY@tc##1{\textcolor[rgb]{0.00,0.50,0.00}{##1}}}
\@namedef{PY@tok@kd}{\let\PY@bf=\textbf\def\PY@tc##1{\textcolor[rgb]{0.00,0.50,0.00}{##1}}}
\@namedef{PY@tok@kn}{\let\PY@bf=\textbf\def\PY@tc##1{\textcolor[rgb]{0.00,0.50,0.00}{##1}}}
\@namedef{PY@tok@kr}{\let\PY@bf=\textbf\def\PY@tc##1{\textcolor[rgb]{0.00,0.50,0.00}{##1}}}
\@namedef{PY@tok@bp}{\def\PY@tc##1{\textcolor[rgb]{0.00,0.50,0.00}{##1}}}
\@namedef{PY@tok@fm}{\def\PY@tc##1{\textcolor[rgb]{0.00,0.00,1.00}{##1}}}
\@namedef{PY@tok@vc}{\def\PY@tc##1{\textcolor[rgb]{0.10,0.09,0.49}{##1}}}
\@namedef{PY@tok@vg}{\def\PY@tc##1{\textcolor[rgb]{0.10,0.09,0.49}{##1}}}
\@namedef{PY@tok@vi}{\def\PY@tc##1{\textcolor[rgb]{0.10,0.09,0.49}{##1}}}
\@namedef{PY@tok@vm}{\def\PY@tc##1{\textcolor[rgb]{0.10,0.09,0.49}{##1}}}
\@namedef{PY@tok@sa}{\def\PY@tc##1{\textcolor[rgb]{0.73,0.13,0.13}{##1}}}
\@namedef{PY@tok@sb}{\def\PY@tc##1{\textcolor[rgb]{0.73,0.13,0.13}{##1}}}
\@namedef{PY@tok@sc}{\def\PY@tc##1{\textcolor[rgb]{0.73,0.13,0.13}{##1}}}
\@namedef{PY@tok@dl}{\def\PY@tc##1{\textcolor[rgb]{0.73,0.13,0.13}{##1}}}
\@namedef{PY@tok@s2}{\def\PY@tc##1{\textcolor[rgb]{0.73,0.13,0.13}{##1}}}
\@namedef{PY@tok@sh}{\def\PY@tc##1{\textcolor[rgb]{0.73,0.13,0.13}{##1}}}
\@namedef{PY@tok@s1}{\def\PY@tc##1{\textcolor[rgb]{0.73,0.13,0.13}{##1}}}
\@namedef{PY@tok@mb}{\def\PY@tc##1{\textcolor[rgb]{0.40,0.40,0.40}{##1}}}
\@namedef{PY@tok@mf}{\def\PY@tc##1{\textcolor[rgb]{0.40,0.40,0.40}{##1}}}
\@namedef{PY@tok@mh}{\def\PY@tc##1{\textcolor[rgb]{0.40,0.40,0.40}{##1}}}
\@namedef{PY@tok@mi}{\def\PY@tc##1{\textcolor[rgb]{0.40,0.40,0.40}{##1}}}
\@namedef{PY@tok@il}{\def\PY@tc##1{\textcolor[rgb]{0.40,0.40,0.40}{##1}}}
\@namedef{PY@tok@mo}{\def\PY@tc##1{\textcolor[rgb]{0.40,0.40,0.40}{##1}}}
\@namedef{PY@tok@ch}{\let\PY@it=\textit\def\PY@tc##1{\textcolor[rgb]{0.24,0.48,0.48}{##1}}}
\@namedef{PY@tok@cm}{\let\PY@it=\textit\def\PY@tc##1{\textcolor[rgb]{0.24,0.48,0.48}{##1}}}
\@namedef{PY@tok@cpf}{\let\PY@it=\textit\def\PY@tc##1{\textcolor[rgb]{0.24,0.48,0.48}{##1}}}
\@namedef{PY@tok@c1}{\let\PY@it=\textit\def\PY@tc##1{\textcolor[rgb]{0.24,0.48,0.48}{##1}}}
\@namedef{PY@tok@cs}{\let\PY@it=\textit\def\PY@tc##1{\textcolor[rgb]{0.24,0.48,0.48}{##1}}}

\def\PYZbs{\char`\\}
\def\PYZus{\char`\_}
\def\PYZob{\char`\{}
\def\PYZcb{\char`\}}
\def\PYZca{\char`\^}
\def\PYZam{\char`\&}
\def\PYZlt{\char`\<}
\def\PYZgt{\char`\>}
\def\PYZsh{\char`\#}
\def\PYZpc{\char`\%}
\def\PYZdl{\char`\$}
\def\PYZhy{\char`\-}
\def\PYZsq{\char`\'}
\def\PYZdq{\char`\"}
\def\PYZti{\char`\~}
% for compatibility with earlier versions
\def\PYZat{@}
\def\PYZlb{[}
\def\PYZrb{]}
\makeatother


    % For linebreaks inside Verbatim environment from package fancyvrb.
    \makeatletter
        \newbox\Wrappedcontinuationbox
        \newbox\Wrappedvisiblespacebox
        \newcommand*\Wrappedvisiblespace {\textcolor{red}{\textvisiblespace}}
        \newcommand*\Wrappedcontinuationsymbol {\textcolor{red}{\llap{\tiny$\m@th\hookrightarrow$}}}
        \newcommand*\Wrappedcontinuationindent {3ex }
        \newcommand*\Wrappedafterbreak {\kern\Wrappedcontinuationindent\copy\Wrappedcontinuationbox}
        % Take advantage of the already applied Pygments mark-up to insert
        % potential linebreaks for TeX processing.
        %        {, <, #, %, $, ' and ": go to next line.
        %        _, }, ^, &, >, - and ~: stay at end of broken line.
        % Use of \textquotesingle for straight quote.
        \newcommand*\Wrappedbreaksatspecials {%
            \def\PYGZus{\discretionary{\char`\_}{\Wrappedafterbreak}{\char`\_}}%
            \def\PYGZob{\discretionary{}{\Wrappedafterbreak\char`\{}{\char`\{}}%
            \def\PYGZcb{\discretionary{\char`\}}{\Wrappedafterbreak}{\char`\}}}%
            \def\PYGZca{\discretionary{\char`\^}{\Wrappedafterbreak}{\char`\^}}%
            \def\PYGZam{\discretionary{\char`\&}{\Wrappedafterbreak}{\char`\&}}%
            \def\PYGZlt{\discretionary{}{\Wrappedafterbreak\char`\<}{\char`\<}}%
            \def\PYGZgt{\discretionary{\char`\>}{\Wrappedafterbreak}{\char`\>}}%
            \def\PYGZsh{\discretionary{}{\Wrappedafterbreak\char`\#}{\char`\#}}%
            \def\PYGZpc{\discretionary{}{\Wrappedafterbreak\char`\%}{\char`\%}}%
            \def\PYGZdl{\discretionary{}{\Wrappedafterbreak\char`\$}{\char`\$}}%
            \def\PYGZhy{\discretionary{\char`\-}{\Wrappedafterbreak}{\char`\-}}%
            \def\PYGZsq{\discretionary{}{\Wrappedafterbreak\textquotesingle}{\textquotesingle}}%
            \def\PYGZdq{\discretionary{}{\Wrappedafterbreak\char`\"}{\char`\"}}%
            \def\PYGZti{\discretionary{\char`\~}{\Wrappedafterbreak}{\char`\~}}%
        }
        % Some characters . , ; ? ! / are not pygmentized.
        % This macro makes them "active" and they will insert potential linebreaks
        \newcommand*\Wrappedbreaksatpunct {%
            \lccode`\~`\.\lowercase{\def~}{\discretionary{\hbox{\char`\.}}{\Wrappedafterbreak}{\hbox{\char`\.}}}%
            \lccode`\~`\,\lowercase{\def~}{\discretionary{\hbox{\char`\,}}{\Wrappedafterbreak}{\hbox{\char`\,}}}%
            \lccode`\~`\;\lowercase{\def~}{\discretionary{\hbox{\char`\;}}{\Wrappedafterbreak}{\hbox{\char`\;}}}%
            \lccode`\~`\:\lowercase{\def~}{\discretionary{\hbox{\char`\:}}{\Wrappedafterbreak}{\hbox{\char`\:}}}%
            \lccode`\~`\?\lowercase{\def~}{\discretionary{\hbox{\char`\?}}{\Wrappedafterbreak}{\hbox{\char`\?}}}%
            \lccode`\~`\!\lowercase{\def~}{\discretionary{\hbox{\char`\!}}{\Wrappedafterbreak}{\hbox{\char`\!}}}%
            \lccode`\~`\/\lowercase{\def~}{\discretionary{\hbox{\char`\/}}{\Wrappedafterbreak}{\hbox{\char`\/}}}%
            \catcode`\.\active
            \catcode`\,\active
            \catcode`\;\active
            \catcode`\:\active
            \catcode`\?\active
            \catcode`\!\active
            \catcode`\/\active
            \lccode`\~`\~
        }
    \makeatother

    \let\OriginalVerbatim=\Verbatim
    \makeatletter
    \renewcommand{\Verbatim}[1][1]{%
        %\parskip\z@skip
        \sbox\Wrappedcontinuationbox {\Wrappedcontinuationsymbol}%
        \sbox\Wrappedvisiblespacebox {\FV@SetupFont\Wrappedvisiblespace}%
        \def\FancyVerbFormatLine ##1{\hsize\linewidth
            \vtop{\raggedright\hyphenpenalty\z@\exhyphenpenalty\z@
                \doublehyphendemerits\z@\finalhyphendemerits\z@
                \strut ##1\strut}%
        }%
        % If the linebreak is at a space, the latter will be displayed as visible
        % space at end of first line, and a continuation symbol starts next line.
        % Stretch/shrink are however usually zero for typewriter font.
        \def\FV@Space {%
            \nobreak\hskip\z@ plus\fontdimen3\font minus\fontdimen4\font
            \discretionary{\copy\Wrappedvisiblespacebox}{\Wrappedafterbreak}
            {\kern\fontdimen2\font}%
        }%

        % Allow breaks at special characters using \PYG... macros.
        \Wrappedbreaksatspecials
        % Breaks at punctuation characters . , ; ? ! and / need catcode=\active
        \OriginalVerbatim[#1,codes*=\Wrappedbreaksatpunct]%
    }
    \makeatother

    % Exact colors from NB
    \definecolor{incolor}{HTML}{303F9F}
    \definecolor{outcolor}{HTML}{D84315}
    \definecolor{cellborder}{HTML}{CFCFCF}
    \definecolor{cellbackground}{HTML}{F7F7F7}

    % prompt
    \makeatletter
    \newcommand{\boxspacing}{\kern\kvtcb@left@rule\kern\kvtcb@boxsep}
    \makeatother
    \newcommand{\prompt}[4]{
        {\ttfamily\llap{{\color{#2}[#3]:\hspace{3pt}#4}}\vspace{-\baselineskip}}
    }
    

    
    % Prevent overflowing lines due to hard-to-break entities
    \sloppy
    % Setup hyperref package
    \hypersetup{
      breaklinks=true,  % so long urls are correctly broken across lines
      colorlinks=true,
      urlcolor=urlcolor,
      linkcolor=linkcolor,
      citecolor=citecolor,
      }
    % Slightly bigger margins than the latex defaults
    
    \geometry{verbose,tmargin=1in,bmargin=1in,lmargin=1in,rmargin=1in}
    
    

\begin{document}
    
    \maketitle
    
    

    
    Probability for Data Science

UC Berkeley, Spring 2026

Ani Adhikari

CC BY-NC-SA 4.0

This content is protected and may not be shared, uploaded, or
distributed.

    \hypertarget{homework-13-due-monday-may-4th-at-5-pm}{%
\section{Homework 13 (Due Monday, May 4th at 5
PM)}\label{homework-13-due-monday-may-4th-at-5-pm}}

    \begin{tcolorbox}[breakable, size=fbox, boxrule=1pt, pad at break*=1mm,colback=cellbackground, colframe=cellborder]
\prompt{In}{incolor}{1}{\boxspacing}
\begin{Verbatim}[commandchars=\\\{\}]
\PY{k+kn}{import}\PY{+w}{ }\PY{n+nn}{warnings}
\PY{n}{warnings}\PY{o}{.}\PY{n}{filterwarnings}\PY{p}{(}\PY{l+s+s1}{\PYZsq{}}\PY{l+s+s1}{ignore}\PY{l+s+s1}{\PYZsq{}}\PY{p}{)}

\PY{k+kn}{from}\PY{+w}{ }\PY{n+nn}{prob140}\PY{+w}{ }\PY{k+kn}{import} \PY{o}{*}
\PY{k+kn}{from}\PY{+w}{ }\PY{n+nn}{datascience}\PY{+w}{ }\PY{k+kn}{import} \PY{o}{*}
\PY{k+kn}{import}\PY{+w}{ }\PY{n+nn}{numpy}\PY{+w}{ }\PY{k}{as}\PY{+w}{ }\PY{n+nn}{np}
\PY{k+kn}{from}\PY{+w}{ }\PY{n+nn}{scipy}\PY{+w}{ }\PY{k+kn}{import} \PY{n}{stats}

\PY{k+kn}{import}\PY{+w}{ }\PY{n+nn}{matplotlib}\PY{n+nn}{.}\PY{n+nn}{pyplot}\PY{+w}{ }\PY{k}{as}\PY{+w}{ }\PY{n+nn}{plt}
\PY{o}{\PYZpc{}}\PY{k}{matplotlib} inline
\PY{k+kn}{import}\PY{+w}{ }\PY{n+nn}{matplotlib}
\PY{n}{matplotlib}\PY{o}{.}\PY{n}{style}\PY{o}{.}\PY{n}{use}\PY{p}{(}\PY{l+s+s1}{\PYZsq{}}\PY{l+s+s1}{fivethirtyeight}\PY{l+s+s1}{\PYZsq{}}\PY{p}{)}
\end{Verbatim}
\end{tcolorbox}

    \hypertarget{instructions}{%
\subsubsection{Instructions}\label{instructions}}

\hypertarget{questions-1-2-and-3-can-be-answered-based-on-tuesdays-lecture-and-past-content-alone.}{%
\paragraph{\texorpdfstring{\textbf{Questions 1, 2, and 3 can be answered
based on Tuesday's lecture (and past content)
alone.}}{Questions 1, 2, and 3 can be answered based on Tuesday's lecture (and past content) alone.}}\label{questions-1-2-and-3-can-be-answered-based-on-tuesdays-lecture-and-past-content-alone.}}

\hypertarget{we-strongly-encourage-you-to-attempt-question-1-relying-on-the-provided-textbook-sections-and-your-own-reasoning.-if-you-find-yourself-stuck-we-recommend-attending-office-hours-or-homework-party-or-posting-your-questions-on-ed.-if-youre-having-trouble-getting-started-on-problems-without-the-use-of-ai-we-encourage-you-to-follow-our-recommendation.}{%
\paragraph{\texorpdfstring{\textbf{We strongly encourage you to attempt
Question 1, relying on the provided textbook sections and your own
reasoning. If you find yourself stuck, we recommend attending office
hours or homework party or posting your questions on Ed. If you're
having trouble getting started on problems without the use of AI, we
encourage you to follow our
recommendation.}}{We strongly encourage you to attempt Question 1, relying on the provided textbook sections and your own reasoning. If you find yourself stuck, we recommend attending office hours or homework party or posting your questions on Ed. If you're having trouble getting started on problems without the use of AI, we encourage you to follow our recommendation.}}\label{we-strongly-encourage-you-to-attempt-question-1-relying-on-the-provided-textbook-sections-and-your-own-reasoning.-if-you-find-yourself-stuck-we-recommend-attending-office-hours-or-homework-party-or-posting-your-questions-on-ed.-if-youre-having-trouble-getting-started-on-problems-without-the-use-of-ai-we-encourage-you-to-follow-our-recommendation.}}

Your homeworks will generally have two components: a written portion and
a portion that also involves code. Written work should be completed on
paper, and coding questions should be done in the notebook. Start the
work for the written portions of each section on a new page. You are
welcome to \(\LaTeX\) your answers to the written portions, but staff
will not be able to assist you with \(\LaTeX\) related issues.

It is your responsibility to ensure that both components of the homework
are submitted completely and properly to Gradescope. \textbf{Make sure
to assign each page of your pdf to the correct question. Refer to the
bottom of the notebook for submission instructions.}

    \hypertarget{how-to-do-your-homework}{%
\subsubsection{How to Do Your Homework}\label{how-to-do-your-homework}}

The point of homework is for you to try your hand at using what you've
learned in class. The steps to follow:

\begin{itemize}
\tightlist
\item
  Go to lecture and sections, and also go over the relevant text
  sections before starting on the homework. This will remind you what
  was covered in class, and the text will typically contain examples not
  covered in lecture. The weekly Study Guide will list what you should
  read.
\item
  Work on some of the practice problems before starting on the homework.
\item
  Attempt the homework problems by yourself with the text, section work,
  and practice materials all at hand. Sometimes the week's lab will help
  as well. The two steps above will help this step go faster and be more
  fruitful.
\item
  At this point, seek help if you need it. Don't ask how to do the
  problem --- ask how to get started, or for a nudge to get you past
  where you are stuck. Always say what you have already tried. That
  helps us help you more effectively.
\item
  For a good measure of your understanding, keep track of the fraction
  of the homework you can do by yourself or with minimal help. It's a
  better measure than your homework score, and only you can measure it.
\end{itemize}

    \hypertarget{rules-for-homework}{%
\subsubsection{Rules for Homework}\label{rules-for-homework}}

\begin{itemize}
\tightlist
\item
  Every answer should contain a calculation or reasoning. For example, a
  calculation such as \((1/3)(0.8) + (2/3)(0.7)\) or
  \texttt{sum({[}(1/3)*0.8,\ (2/3)*0.7{]})}is fine without further
  explanation or simplification. If we want you to simplify, we'll ask
  you to. But just \({5 \choose 2}\) by itself is not fine; write ``we
  want any 2 out of the 5 frogs and they can appear in any order'' or
  whatever reasoning you used. Reasoning can be brief and abbreviated,
  e.g.~``product rule'' or ``not mutually exclusive.''
\item
  You may consult others (see ``How to Do Your Homework'' above) but you
  must write up your own answers using your own words, notation, and
  sequence of steps.
\item
  We'll be using Gradescope. You must submit the homework according to
  the instructions at the end of homework set.
\end{itemize}

\hypertarget{we-will-not-grade-assignments-which-do-not-have-pages-correctly-selected-for-each-question.}{%
\subsection{\texorpdfstring{\textbf{We will not grade assignments which
do not have pages correctly selected for each
question.}}{We will not grade assignments which do not have pages correctly selected for each question.}}\label{we-will-not-grade-assignments-which-do-not-have-pages-correctly-selected-for-each-question.}}

    \newpage

    \hypertarget{predicting-scores}{%
\subsection{1. Predicting Scores}\label{predicting-scores}}

{[}Your answers to this question should be decimal values or equations
with numerical coefficients. For the arithmetic, you are welcome use the
code cell at the end of the question. It's just there for your
convenience -- we won't read it.{]}

Grades in a class are based on a linear combination of a final exam
(worth 50\% of the grade), a midterm (worth 30\%), and homework (worth
20\%). Let the random vector \([F ~~ M ~~ H]^T\) consist of the final,
midterm, and homework scores of a randomly picked student.

Suppose the mean vector of \([F ~~ M ~~ H]^T\) is \([60 ~~ 55 ~~ 80]^T\)
and the covariance matrix is

\[
\begin{bmatrix}
121 & 80 & 10 \\
80 & 144 & 15 \\
10 & 15 & 9
\end{bmatrix}
\]

Suppose the distribution of \([F ~~ M ~~ H]^T\) is multivariate normal.

\textbf{a)} Find the distribution of the student's overall score
\(S = 0.5F + 0.3M + 0.2H\). Read through
\href{https://data140.org/textbook/content/chapter-23/random-vectors/}{Ch.
23.1} (Random Vectors) and
\href{https://data140.org/textbook/content/chapter-23/multivariate-normal-vectors/}{Ch.
23.2} (Multivariate Normal Vectors) to get started.

Let $A$ be the weights. 
\[
	S = \underbrace{\begin{bmatrix}
	0.5 &0.3 &0.2	
\end{bmatrix}}_{A} \begin{bmatrix}
F \\ M\\H
\end{bmatrix}
\]

\begin{align*}
E[S] &=  E[0.5F + 0.3M + 0.2H] \\
&= E[0.5F]  + E[0.3M] + E[0.2H]\\ 
&= 0.5(60) + 0.3(55) + 0.2(80) \\
&= 62.5 \\
\end{align*}


\begin{align*}
	\Sigma_{S} &= A \sigma A^{T} \\
	&= 71.37 \\
\end{align*}

Therefore $S \sim Normal(62.5, 71.37)$ 

\textbf{b)} The instructor wonders whether the final exam score \(F\)
can just be predicted by a linear function of the random variable
\(X = 0.3M + 0.2H\). Explain whether the least squares predictor of
\(F\) based on \(X\) is linear and why. Then find the least squares
predictor of \(F\) based on \(X\). Relevant background can be found in
\href{https://data140.org/textbook/content/chapter-24/linear-least-squares/\#slope-and-intercept-of-the-regression-line}{Ch.
24.1.2} (Slope and Intercept of the Regression Line) and
\href{https://data140.org/textbook/content/chapter-24/regression-and-bivariate-normal/}{Ch.
24.3} (Regression and the Bivariate Normal).



Yes, the final exam score can be predicted by a linear function. $X$ is a linear combination of a multivariate normal distribution. Therefore, $\begin{bmatrix}
	F &X
\end{bmatrix}^{T}$ is a bivariate normal distribution as well. 



We showed in lecture that the least squares predictor of a bivariate distribution is linear. Therefore, since $\begin{bmatrix}
	F &X	
\end{bmatrix}^{T}$ are bivariate normal, then the least squares predictor is linear. 


$X \sim Normal(32.5, 15.12)$  

(This used the code cell to calculate the values) 

The least squares predictor of $F$ based on $X$ is thus:


\begin{align*}
	E[F \mid X] &= E[\sigma_{F}F_{su} + \mu_{F} \mid X_{su}] \\
	&= \sigma_{F}E[F_{su} \mid X_{su}] + \mu _{F} \\
	&= \sigma_{F} \rho \left( \frac{X-\mu _{X}}{\sigma_{X}} \right)   + \mu _{F}\\ 
	&= \rho \frac{\sigma_{F}}{\sigma_{X}} \left( X-\mu _{X} \right) +\mu _{F} \\
\end{align*} 

\begin{align*}
	\rho \frac{\sigma_{F}}{\sigma_{X}}  &= \frac{\Cov(X, F)}{\Var(X)} \\
	&= \frac{\Cov(0.3M, F) + \Cov(0.2H, F)}{\Var(0.3M + 0.2H)} \\
	&= \frac{0.3 \Cov(M,F) + .2\Cov(H,F)}{0.3^{2} \Var(M) + 0.2^{2}\Var(H) + 2\Cov(0.3M, 0.2H)} \\
	&= \frac{.3 (80) + .2 (10)}{0.3^{2} \cdot  144 + 0.2^{2} \cdot 2 \cdot 0.3 \cdot .2 \Cov(M, H) } \\
	&= \frac{.3(80) + .2 (10)}{0.3^{2} \cdot 144 + 0.2^{2} \cdot 9 + .12 \cdot 15} \\
	&= \frac{26}{15.12} \\
	&= 1.72 \\
\end{align*}

\[
E[F \mid X] = 1.72 X  + 4.1
\]

\textbf{c)} Find the root mean squared error of the predictor in Part
\textbf{b}. Utilize the equation in
\href{https://data140.org/textbook/content/chapter-24/regression-and-bivariate-normal/\#prediction-error}{Ch.
24.3.1} (Prediction Error) to find the MSE, but remeber that this
equation is for when X and Y are standard bivariate normal. This is not
the case for this question, so the setup will be slightly different.
Additionally, RMSE is the square root of the MSE.



\begin{align*}
	MSE(E[F \mid X]) &= E[ \left( F - E[F \mid X] \right) ^{2}] \\
	&= E[E[\left( F - E[F \mid X] \right) ^{2} \mid X]] \\
	&= E[\Var(F \mid X)] \\
\end{align*}

\begin{align*}
	\Var(F \mid X) &= \Var(\sigma_{X} \left( \rho \left( \frac{f-\mu _{F}}{\sigma_{F}} \right)  + \sqrt{1-\rho^{2}} Z \right)  + \mu _{X}) \\
		       &= \Var(\sigma_{X} \left( \sqrt{1-\rho^{2}} Z \right))  \\
		       &= \sigma_{X}^{2} (1-\rho^{2}) \\
\end{align*}


\[
E[\Var(F \mid X)] = \sigma_{X}^{2} (1-\rho^{2})
\]
RMSE




    \begin{tcolorbox}[breakable, size=fbox, boxrule=1pt, pad at break*=1mm,colback=cellbackground, colframe=cellborder]
\prompt{In}{incolor}{ }{\boxspacing}
\begin{Verbatim}[commandchars=\\\{\}]
\PY{n}{weight} \PY{o}{=} \PY{n}{np}\PY{o}{.}\PY{n}{matrix}\PY{p}{(}\PY{p}{[}\PY{l+m+mf}{0.5}\PY{p}{,} \PY{l+m+mf}{0.3}\PY{p}{,} \PY{l+m+mf}{0.2}\PY{p}{]}\PY{p}{)}
\PY{n}{mu} \PY{o}{=} \PY{n}{np}\PY{o}{.}\PY{n}{matrix}\PY{p}{(}\PY{p}{[}\PY{l+m+mi}{60}\PY{p}{,} \PY{l+m+mi}{55}\PY{p}{,} \PY{l+m+mi}{80}\PY{p}{]}\PY{p}{)}
\PY{n}{cov} \PY{o}{=} \PY{n}{np}\PY{o}{.}\PY{n}{matrix}\PY{p}{(}\PY{p}{[}\PY{p}{[}\PY{l+m+mi}{121}\PY{p}{,} \PY{l+m+mi}{80}\PY{p}{,} \PY{l+m+mi}{10}\PY{p}{]}\PY{p}{,} 
      \PY{p}{[}\PY{l+m+mi}{80}\PY{p}{,} \PY{l+m+mi}{144}\PY{p}{,} \PY{l+m+mi}{15}\PY{p}{]}\PY{p}{,}
      \PY{p}{[}\PY{l+m+mi}{10}\PY{p}{,} \PY{l+m+mi}{15}\PY{p}{,} \PY{l+m+mi}{9}\PY{p}{]}\PY{p}{]}\PY{p}{)}

\PY{o}{.}\PY{o}{.}\PY{o}{.}
\end{Verbatim}
\end{tcolorbox}

    \newpage

    \hypertarget{normal-sample-mean-and-sample-variance-part-1}{%
\subsection{2. Normal Sample Mean and Sample Variance, Part
1}\label{normal-sample-mean-and-sample-variance-part-1}}

Let \(X_1, X_2, \ldots, X_n\) be i.i.d. with mean \(\mu\) and variance
\(\sigma^2\). Let

\[
\bar{X} ~ = ~ \frac{1}{n} \sum_{i=1}^n X_i
\]

denote the sample mean. In
\href{http://prob140.datahub.berkeley.edu/hub/user-redirect/git-pull?repo=https://github.com/prob140/materials-sp26\&branch=main\&subPath=hw/Homework_07.ipynb}{Homework
7}, you constructed a random variable

\[
S^2 = \frac{1}{n - 1} \sum_{i=1}^n (X_i - \bar{X})^2
\]

called the sample variance. Before proceeding, please review Homework 7,
Question 3 in its entirety.

\textbf{a)} For \(1 \le i \le n\) let \(D_i = X_i - \bar{X}\). Find
\(Cov(D_i, \bar{X})\).

\begin{align*}
	\Cov(D_{i}, \bar X) &= \Cov(X_{i}, \bar X) + \Var(\bar X) \\
			    &= \frac{\sigma^{2}}{n} - \frac{\sigma^{2}}{n}, \text{by HW 7} \\
			    &= 0 \\
\end{align*}




\textbf{b)} Now assume in addition that \(X_1, X_2, \ldots, X_n\) are
i.i.d. normal \((\mu, \sigma^2)\). What is the joint distribution of
\(\bar{X}, D_1, D_2, \ldots, D_{n-1}\)? Explain why \(D_n\) isn't on the
list.


$D_{n}$ isn't on the list because it is fully explained by the other variables. In other words, it is 


\textbf{c)} True or false (justify your answer): The sample mean and
sample variance of an i.i.d. normal sample are independent of each
other.

    \begin{tcolorbox}[breakable, size=fbox, boxrule=1pt, pad at break*=1mm,colback=cellbackground, colframe=cellborder]
\prompt{In}{incolor}{ }{\boxspacing}
\begin{Verbatim}[commandchars=\\\{\}]

\end{Verbatim}
\end{tcolorbox}

    \newpage

    \hypertarget{normal-sample-mean-and-sample-variance-part-2}{%
\subsection{3. Normal Sample Mean and Sample Variance, Part
2}\label{normal-sample-mean-and-sample-variance-part-2}}

\textbf{a)} Let \(R\) have the chi-squared distribution with \(n\)
degrees of freedom. What is the mgf of \(R\)?


$R \sim gamma(\frac{n}{2}, \frac{1}{2})$ 


\begin{align*}
	M(t) = \left( \frac{\frac{1}{2}}{\frac{1}{2} - t} \right) ^{\frac{n}{2}} 
\end{align*}


\textbf{b)} For \(R\) as in Part (a), suppose \(R = V + W\) where \(V\)
and \(W\) are independent and \(V\) has the chi-squared distribution
with \(m < n\) degrees of freedom. Can you identify the distribution of
\(W\)? Justify your answer.

We can express the MGF of R in terms of V and W. We can equate the exponents to make the equality true. 
\begin{align*}
	M_{R}(t) &=  M_{V}(t) M_{W}(t) \\
	\left( \frac{\frac{1}{2}}{\frac{1}{2}-t} \right) ^{\frac{n}{2}} &=  \left( \frac{\frac{1}{2}}{\frac{1}{2}-t} \right) ^{\frac{m}{2}} \cdot  \left( \frac{\frac{1}{2}}{\frac{1}{2}-t} \right) ^{\frac{n-m}{2}}  \\
\end{align*} 


Therefore, $W \sim gamma(\frac{n-m}{2}, \frac{1}{2})$


\textbf{c)} Let \(X_1, X_2, \ldots , X_n\) be any sequence of random
variables and let \(\bar{X} = \frac{1}{n}\sum_{i=1}^n X_i\). Let
\(\alpha\) be any constant. Prove the \emph{sum of squares
decomposition}

\[
\sum_{i=1}^n (X_i - \alpha)^2 ~=~ \sum_{i=1}^n (X_i - \bar{X})^2 ~+~ n(\bar{X} - \alpha)^2
\]


\begin{align*}
\sum_{i=1}^{n} (X_{i} - \alpha)^{2} &= \sum_{i=1}^{n} (X_{i} - \bar X + \bar X - \alpha)^{2} \\
				    &= \sum_{i=1}^{n} \left[ (X_{i} - \bar X) + (\bar X - \alpha) \right]^{2} \\
				    &= \sum_{i=1}^{n} (X_{i} - \bar X)^{2} + \sum_{i=1}^{n} (\bar X - \alpha)^{2} + \sum_{i=1}^{n} 2(X_{i} - \bar X)(\bar X - \alpha) \\
		    &= \sum_{i=1}^{n} (X_{i} - \bar X)^{2} + n (\bar X - \alpha)^{2} + 2(\bar X - \alpha) \sum_{i=1}^{n} (X_{i}- \bar X) \\
		    &= \sum_{i=1}^{n}(X_{i} - \bar X)^{2} + n(\bar X - \alpha)^{2} + 2(\bar X - \alpha) \left[ \sum_{i=1}^{n} X_{i} - n \bar X \right]\\
		    &= \sum_{i=1}^{n}(X_{i} - \bar X)^{2} + n(\bar X - \alpha)^{2}+ 2(\bar X - \alpha) \left[ \sum_{i=1}^{n}X_{i} - n \frac{1}{n}\sum_{i=1}^{n} X_{i}\right]  \\
		    &= \sum_{i=1}^{n}(X_{i} - \bar X)^{2} + n(\bar X - \alpha)^{2}   \\
\end{align*}








\textbf{d)} Now let \(X_1, X_2, \ldots , X_n\) be i.i.d. normal with
mean \(\mu\) and variance \(\sigma^2 > 0\). Let \(S^2\) be the ``sample
variance'' defined by

\[
S^2 ~=~ \frac{1}{n-1} \sum_{i=1}^n (X_i - \bar{X})^2
\]

Find a constant \(c\) such that \(cS^2\) has a chi-squared distribution.
Provide the degrees of freedom.

{[}Use Parts (b) and (c) as well as the result of the previous
exercise.{]}

\begin{align*}
	\sum_{i=1}^{n}(X_{i} - \mu)^{2} &= \sum_{i=1}^{n}(X_{i}- \bar X)^{2} + n(\bar X - \mu )^{2} \\
	\sum_{i=1}^{n} (X_{i} - \mu )^{2} &=  (n-1)S^{2} + n(\bar X - \mu )^{2} \\
	\sum_{i=1}^{n} \left( \frac{X_{i} - \mu }{\sigma} \right)^{2} &= \frac{n-1}{\sigma^{2}}S^{2} + \left(\frac{n(\bar X - \mu )}{\sigma}\right)^{2} \\
	\sum_{i=1}^{n}Z^{2}  &= \frac{n-1}{\sigma^{2}} S^{2} + \left( \frac{\bar X - \mu }{\sigma / n} \right) ^{2} \\
	\sum_{i=1}^{n} Z^{2} &= \frac{n-1}{\sigma^{2}} S^{2} + Z^{2} \\
	\underbrace{	\sum_{i=1}^{n} (Z^{2})}_{gamma(\frac{n}{2}, \frac{1}{2})} - \underbrace{Z^{2}}_{gamma(\frac{1}{2}, \frac{1}{2})} &= \frac{n-1}{\sigma^{2}}S^{2} \\
\end{align*}


$c = \frac{n-1}{\sigma^{2}} S^{2}$ and The LHS is gamma($\frac{n-1}{2}, \frac{1}{2})$, therefore the degrees of freedom is $n-1$ 


    \newpage

    \hypertarget{multiple-regression-residuals}{%
\subsection{4. Multiple Regression:
Residuals}\label{multiple-regression-residuals}}

This exercise assumes the multiple regression model of
\href{https://data140.org/textbook/content/chapter-25/multiple-regression-1/}{Section
25.4} of the textbook and uses the same notation as in that section.

\textbf{a)} Find the distribution of the parameter estimate vector
\(\hat{\boldsymbol{\beta}}\). Derive the answer; don't just quote a
result from lecture.

From lecture we established that: 

\[
Y = X \beta + \epsilon
\]
\[
Y = X \hat \beta + e
\]

Which means that the true $Y$ is our estimated signals and some estimated noise. We also established that $Y$ is multivariate normal with mean vector $X \beta$, and covariance matrix $\sigma^{2}I_{n}$, because we assume errors are multivariate normal.

We also established that if $\mathbb{X}$ is full rank, then $\hat \beta = (X^{T}X)^{-1}X^{T}Y$. $(X^{T}X)^{-1}X^{T}$ is a linear transformation matrix, and we can call this $A$. 

Then since we know that $Y \sim $ Multivariatenormal($X\beta, \sigma^{2}I_{n}$), then under the linear transformation rules for Multivariatenormals, $\hat \beta \sim $  multivariatenormal($A X\beta, \sigma^{2} AI_{n}A^{T} $ )

Expanding it out: 

\begin{align*}
AX\beta &= (X^{T}X)^{-1} X^{T}X \beta  \\
&= \beta \\
\end{align*}


\begin{align*}
	\sigma^{2}AI_{n}A^{T} &= \sigma A A^{T} \\
	&= \sigma^{2} (X^{T}X)^{-1}X^{T} X (X^{T}X)^{-1^{T}} \\
	&= \sigma^{2} (X^{T}X)^{-1} \\
\end{align*}
The distribution of $\hat \beta$ is MultivariateNormal($\beta, \sigma^{2}(X^{T}X)^{-1}$) 

\textbf{b)} The regression estimate \(\hat{\mathbf{Y}}\) can be written
as \(\mathbf{HY}\) for an \(n \times n\) matrix \(\mathbf{H}\). This
matrix is called the \emph{hat matrix}, probably because it `'puts the
hat on \(\mathbf{Y}\).'' Write \(\mathbf{H}\) in terms of
\(\mathbf{X}\). Is \(\mathbf{H}\) symmetric?

$$
	\hat Y = HY
$$
\begin{align*}
	\hat Y &= X \hat \beta \\
	&= X (X^{T}X)^{-1}X^{T}Y \\
\end{align*}


\[
H = X(X^{T}X)^{-1}X^{T}
\]

\begin{align*}
	H^{T} &= \left( X(X^{T}X)^{-1}X^{T} \right)^{T}  \\
	      &= X^{T^{T}} (X^{T}X)^{-1^{T}}X^{T} \\
&= X^{T} (X^{T}X)^{-1}X^{T} \\
\end{align*}

Therefore $H = H^{T}$, making it symmetric
\textbf{c)} Show that \(\mathbf{H}\) is idempotent for \(\mathbf{H}\)
defined in Part (b). (If you haven't seen that term before, look it up.)

Idempotent is if $H$ is applied multiple times, it is the same operation. In simpler terms, $H^{n} = H$ 


\begin{align*}
	H^{2} &= X\underbrace{(X^{T}X)^{-1}X^{T} X}_{\mathbb{I}}(X^{T}X)^{-1} X^{T} \\
	&=  X (X^{T}X)^{-1}X^{T}\\
\end{align*}



\textbf{d)} Find the distribution of the residual vector \(\mathbf{e}\).

\[
e = Y - \hat Y
\]

\begin{align*}
e &= Y - HY \\
&= (I-H)Y \\
\end{align*}

$e$ is multivariatenormal($(I-H) \mu _{Y}, (I-H) \Sigma_{Y} (I-H)^{T}$  )



\textbf{e)} Show that the covariance matrix of \(\mathbf{e}\) is
\(\sigma^2(\mathbf{I} - \mathbf{H})\).



\begin{align*}
	(I-H)\Sigma_Y(I-H)^{T} &= (I-H) \sigma^{2}I_{n} (I-H)^{T} \\
	&= \sigma^{2} (I-H) (I-H)^{T} \\
	&= \sigma^{2} (I - H) ( I^{T} - H^{T}) \\
	&= \sigma^{2}(I-H)(I - H), \text{by symmetric matricies} \\
	&= \sigma^{2}( I^{2} - IH - HI + H^{2}) \\
	&= \sigma^{2} (I^{2} - H - H + H^{2}) \\
	&= \sigma^{2} (I - 2H + H) , \text{by idempotent H}\\
	&= \sigma^{2} (I - H) \\
\end{align*}







    \newpage

    \hypertarget{submission-instructions}{%
\subsection{Submission Instructions}\label{submission-instructions}}

Many assignments throughout the course will have a written portion and a
code portion. Please follow the directions below to properly submit both
portions.

\hypertarget{written-portion}{%
\subsubsection{Written Portion}\label{written-portion}}

\begin{itemize}
\tightlist
\item
  Scan all the pages into a PDF. You can use any scanner or a phone
  using applications such as CamScanner. Please \textbf{DO NOT} simply
  take pictures using your phone.
\item
  Please start a new page for each question. If you have already written
  multiple questions on the same page, you can crop the image in
  CamScanner or fold your page over (the old-fashioned way). This helps
  expedite grading.
\item
  It is your responsibility to check that all the work on all the
  scanned pages is legible.
\item
  If you used \(\LaTeX\) to do the written portions, you do not need to
  do any scanning; you can just download the whole notebook as a PDF via
  LaTeX.
\end{itemize}

\hypertarget{code-portion}{%
\subsubsection{Code Portion}\label{code-portion}}

\begin{itemize}
\tightlist
\item
  Save your notebook using
  \texttt{File\ \textgreater{}\ Save\ and\ Checkpoint}.
\item
  Generate a PDF file using
  \texttt{File\ \textgreater{}\ Download\ As\ \textgreater{}\ PDF\ via\ LaTeX}.
  This might take a few seconds and will automatically download a PDF
  version of this notebook.

  \begin{itemize}
  \tightlist
  \item
    If you have issues, please post a follow-up on the general Homework
    13 Ed thread.
  \end{itemize}
\end{itemize}

\hypertarget{submitting}{%
\subsubsection{Submitting}\label{submitting}}

\begin{itemize}
\tightlist
\item
  Combine the PDFs from the written and code portions into one PDF.
  \href{https://smallpdf.com/merge-pdf}{Here} is a useful tool for doing
  so.
\item
  Submit the assignment to Homework 13 on Gradescope.
\item
  \textbf{Make sure to assign each page of your pdf to the correct
  question.}
\item
  \textbf{It is your responsibility to verify that all of your work
  shows up in your final PDF submission.}
\end{itemize}

If you are having difficulties scanning, uploading, or submitting your
work, please read the
\href{https://edstem.org/us/courses/94054/discussion/7533569}{Ed Thread}
on this topic and post a follow-up on the general Homework 13 Ed thread.


    % Add a bibliography block to the postdoc
    
    
    
\end{document}
